La evaluación del desempeño ambiental de la bioconversión con \textit{Hermetia illucens} requiere analizar sus emisiones gaseosas en comparación con tecnologías convencionales de tratamiento biológico, especialmente el compostaje. Aunque ambos procesos se basan en rutas aeróbicas de transformación de materia orgánica, la bioconversión presenta una dinámica distintiva. Esta es determinada por la interacción entre el metabolismo larval, la actividad microbiana y los cambios acelerados en las propiedades fisicoquímicas del sustrato \parencite{rossiEstimatingDynamicsGreenhouse2024a,eriksenDynamicModellingFeed2022,parodiBioconversionEfficienciesGreenhouse2020a}. 

Desde una perspectiva ambiental, esta diferencia es crucial, ya que el perfil de emisiones no depende únicamente del tipo de residuo, sino también de la velocidad de procesamiento, la distribución de oxígeno y la partición del carbono y del nitrógeno entre la biomasa, el residuo transformado (\textit{frass}) y los flujos gaseosos \parencite{parodiBioconversionEfficienciesGreenhouse2020a,pangReducingGreenhouseGas2020,ermolaevGreenhouseGasEmissions2019}.

\subsection{Dinámica de generación de emisiones en el sistema larva--sustrato}

Durante la bioconversión, el \ce{CO2} se genera principalmente por la respiración aeróbica de las larvas y de la microbiota asociada. Estudios clásicos estiman que una fracción importante del carbono ingerido se libera como \ce{CO2}, lo que convierte a este gas en un indicador central de la intensidad metabólica del sistema \parencite{dienerConversionOrganicMaterial2011,eriksenDynamicModellingFeed2022}. No obstante, su trayectoria temporal depende de variables críticas como la composición del sustrato, el contenido de humedad, la densidad larval, la temperatura y la disponibilidad de oxígeno \parencite{rossiEstimatingDynamicsGreenhouse2024a,bekkerImpactSubstrateMoisture2021a}.

A diferencia del compostaje, donde la degradación es exclusivamente microbiana y se desarrolla en escalas de tiempo prolongadas, la bioconversión incorpora la acción simultánea de fragmentación mecánica larval y consumo rápido de materia orgánica. Esta interacción acelera la reducción de masa y modifica con rapidez variables como el pH, la porosidad y la disponibilidad de nutrientes. Tales dinámicas generan microambientes que pueden favorecer o restringir rutas metabólicas específicas a lo largo del proceso \parencite{surendraBioconversionOrganicWastes2016,bermudez-serranoComprehensiveUtilizationBlack2023,rossiEstimatingDynamicsGreenhouse2024a}. Así, las emisiones gaseosas dejan de ser un subproducto uniforme para convertirse en la expresión dinámica del estado metabólico del sistema.



\subsection{Perfil comparativo: \textit{Hermetia illucens} frente al compostaje tradicional}

La literatura sugiere que la bioconversión con \textit{Hermetia illucens} presenta un perfil de emisiones más favorable que el compostaje tradicional, particularmente respecto a gases de alto potencial de calentamiento global como el metano (\ce{CH4}) y el óxido nitroso (\ce{N2O}) \parencite{pangReducingGreenhouseGas2020,ermolaevGreenhouseGasEmissions2019,parodiBlackSoldierFly2021a,mertenatBlackSoldierFly2019}. 

En el compostaje, la formación de zonas anaeróbicas núcleos de la pila suele generar emisiones apreciables de \ce{CH4} si el control de aireación es insuficiente \parencite{nordahlGreenhouseGasAir2023a}. En sistemas con \textit{Hermetia illucens}, la actividad de excavación y movimiento larval modifica la estructura del medio, mejorando la porosidad y reduciendo la permanencia del residuo en condiciones anaeróbicas propicias para la metanogénesis \parencite{pangReducingGreenhouseGas2020,parodiBlackSoldierFly2021a}. 

De forma similar, las pérdidas gaseosas de nitrógeno como \ce{N2O} tienden a disminuir cuando el nitrógeno del sustrato es incorporado eficientemente a la biomasa larval \parencite{parodiBioconversionEfficienciesGreenhouse2020a,parodiBlackSoldierFly2021a}. En este contexto, el \ce{CO2} permanece como el principal flujo gaseoso en ambos procesos; sin embargo, en la bioconversión este gas adquiere un valor analítico particular al reflejar de forma continua la actividad biológica agregada, permitiendo su uso como variable de seguimiento del estado del proceso \parencite{eriksenDynamicModellingFeed2022,rossiEstimatingDynamicsGreenhouse2024a}.

\subsection{Evidencia cuantitativa e incertidumbre técnica}

Pese al interés científico, la evidencia cuantitativa sigue siendo fragmentaria y metodológicamente heterogénea. Se han reportado emisiones de \ce{CO2} entre 7.76 y 11.88\,g por gramo de biomasa larval seca \parencite{rossiEstimatingDynamicsGreenhouse2024a}, así como pérdidas gaseosas de carbono cercanas al 24\,\% de la entrada dietaria en balances de masa \parencite{parodiBioconversionEfficienciesGreenhouse2020a}. 

Esta variabilidad evidencia una brecha crítica entre el potencial ambiental atribuido a la tecnología y la capacidad real de monitorear y predecir sus emisiones bajo criterios consistentes. La validación de la bioconversión como estrategia de mitigación verificable depende de métodos robustos para la medición continua de flujos gaseosos \parencite{nordahlGreenhouseGasAir2023a,rossiEstimatingDynamicsGreenhouse2024a}. En este escenario, el \ce{CO2} se consolida no solo como una emisión, sino como una variable observacional clave para el desarrollo de herramientas de monitoreo predictivo en tiempo real.